\documentclass[aspectratio=169,11pt]{beamer}

\usepackage{fontspec}
\usepackage{graphicx}
\usepackage{booktabs}
\usepackage{tikz}
\usepackage{ragged2e}
\usepackage{dsfont}

\setsansfont{Arial}
\setbeamertemplate{navigation symbols}{}
\setbeamertemplate{footline}{}
\setbeamercolor{normal text}{fg=black,bg=white}
\setbeamercolor{frametitle}{fg=black,bg=white}
\definecolor{deckblue}{HTML}{234E9A}
\definecolor{decknavy}{HTML}{162B4E}
\definecolor{deckgray}{HTML}{68717F}

\setbeamerfont{title}{series=\bfseries,size=\fontsize{28}{32}}
\setbeamerfont{subtitle}{size=\fontsize{15}{19}}
\setbeamerfont{section title}{series=\bfseries,size=\fontsize{27}{32}}
\setbeamertemplate{itemize item}{\raise0.15ex\hbox{\large\textbullet}}
\setlength{\leftmargini}{1.3em}
\setlength{\leftmarginii}{1.1em}

\title{Heaviside Composite\\Optimization Problems}
\subtitle{}
\author{}
\date{}

\begin{document}

\begin{frame}[plain]
  \begin{tikzpicture}[remember picture,overlay]
    \fill[decknavy] (current page.south west) rectangle (current page.north east);
    \fill[deckblue] (current page.south west) rectangle
      ([xshift=0.18cm]current page.north west);
  \end{tikzpicture}
  \vspace{1.65cm}
  {\color{white}\usebeamerfont{title}\inserttitle\par}
  \vfill
\end{frame}

\section{Introduction}

\begin{frame}[plain]
  \begin{tikzpicture}[remember picture,overlay]
    \fill[deckblue] (current page.south west) rectangle
      ([xshift=0.18cm]current page.north west);
    \draw[deckgray!35,line width=0.5pt]
      ([xshift=1.25cm,yshift=-0.4cm]current page.west) --
      ([xshift=-1.25cm,yshift=-0.4cm]current page.east);
  \end{tikzpicture}
  \vfill
  {\color{deckgray}\large SECTION 1\par}
  \vspace{0.25cm}
  {\color{decknavy}\usebeamerfont{section title}Introduction\par}
  \vfill
\end{frame}

\begin{frame}[plain,t]
  \vspace{0.02cm}
  {\color{decknavy}\resizebox{\textwidth}{!}{%
    A bird's eye-view of post World-War-II optimization till now:}\par}

  \vspace{0.04cm}
  \begin{itemize}
    \setlength{\itemsep}{0.035cm}
    \setlength{\parskip}{0pt}
    \fontsize{13.0}{15.2}\selectfont
    \item From linear programming to nonlinear and integer programming.
    \item From nonlinear programming to stochastic programming and modern-day convex programming.
    \item From differentiability to nondifferentiability.
    \item From nondifferentiability to semicontinuity.
    \item From convexity to nonconvexity.
  \end{itemize}

  \vspace{-0.02cm}
  {\color{deckblue}\fontsize{10.8}{13.1}\selectfont
    Y. Cui and J. S. Pang, \textit{Modern Nonconvex Nondifferentiable
    Optimization}. SIAM, MOS--SIAM Series on Optimization
    (December 2021).\par}

  \vspace{0.10cm}
  {\color{decknavy}\rule{0.48\textwidth}{0.45pt}\par}
  \vspace{0.08cm}

  {\color{deckblue}\fontsize{12.0}{14.4}\selectfont
    Apart from further developments of the present, what is the future??\par}
  \vspace{0.06cm}
  {\color{decknavy}\fontsize{11.5}{13.8}\selectfont
    From semicontinuity to lack thereof $\Rightarrow$ discontinuous optimization.\par
    \vspace{0.08cm}
    From closedness of constraints to non-closedness.\par}
\end{frame}

\begin{frame}[plain,t]
  \vspace{0.02cm}
  {\color{decknavy}\bfseries\resizebox{\textwidth}{!}{%
    The General Heaviside Composite Optimization Problem}\par}
  {\color{decknavy}\fontsize{14.6}{17.0}\selectfont (G-HSCOP)\par}

  \vspace{0.10cm}
  \begingroup
  \setlength{\fboxsep}{0.20cm}
  \setlength{\fboxrule}{0.45pt}
  \color{deckgray}
  \noindent\fbox{%
    \begin{minipage}{0.94\textwidth}
      \color{decknavy}
      \fontsize{11.2}{13.3}\selectfont
      \vspace{-0.08cm}
      \noindent\resizebox{0.93\linewidth}{!}{$\displaystyle
        \underset{x\in P}{\operatorname{maximize}}\quad
        \Phi(x)\triangleq c(x)+\sum_{k=1}^{K_0}\psi_{0k}(x)
        \underbrace{\mathds{1}_{[0,\infty)}\!\left(\phi_{0k}(x)\right)}_
          {\text{composite indicator fnc.}}
      $}\par
      \vspace{0.08cm}
      \noindent\resizebox{0.96\linewidth}{!}{$\displaystyle
        \operatorname{subject\ to}\quad
        A_{i\bullet}x+\sum_{k=1}^{K_i}\psi_{ik}(x)
        \mathds{1}_{[0,\infty)}\!\left(\phi_{ik}(x)\right)
        \geq \eta_i,\qquad i=1,\ldots,I.
      $}\par
      \vspace{-0.06cm}
    \end{minipage}%
  }
  \endgroup

  \vspace{0.10cm}
  \begin{itemize}
    \setlength{\itemsep}{0.08cm}
    \setlength{\parskip}{0pt}
    \fontsize{10.25}{12.1}\selectfont
    \item $P$ is a \textcolor{red!72!black}{friendly} set, $c$ is a
      \textcolor{red!72!black}{friendly} function, e.g., affine or quadratic;
      \textcolor{red!72!black}{friendly} $=$ amenable to well-known treatments,
      and $A_{i\bullet}$ is the $i$-th row of a matrix
      $A\in\mathbb{R}^{I\times n}$.
    \item $\psi_{ik}$ and $\phi_{ik}$ are functions whose properties highlight
      the challenges;
    \item most importantly, $\mathds{1}_{[0,\infty)}$ is the ``closed''
      Heaviside function given by
  \end{itemize}

  \vspace{-0.17cm}
  {\color{decknavy}\fontsize{10.5}{12.3}\selectfont
  \[
    \mathds{1}_{[0,\infty)}(s)\triangleq
    \begin{cases}
      1, & \text{if }s\in[0,\infty),\\
      0, & \text{otherwise},
    \end{cases}
    \qquad\text{(upper semi-continuous).}
  \]
  }
\end{frame}

\end{document}
